\documentclass[
	fontsize=12pt,
	paper=a4,
	numbers=noenddot,
	parskip=half,
	openany
]{scrbook}
\usepackage{tcolorbox}

\include{../../../for_latex/preamble}
% ##################################################
% Gemeinsames Setup fuer die Probeklausuren
% (Java-Listings, Hinweis-/Beispiel-Boxen, Struktogramm-Stile in TikZ)
% wird nach \include{../for_latex/preamble} per \input eingebunden
% ##################################################

% --- Java Code-Highlighting ---
\definecolor{codegray}{rgb}{0.5,0.5,0.5}
\definecolor{codepurple}{rgb}{0.58,0,0.82}
\definecolor{backcolour}{rgb}{0.95,0.95,0.92}
\definecolor{keywordcolor}{rgb}{0.8,0.47,0.13}

\lstdefinestyle{javastyle}{
	backgroundcolor=\color{backcolour},
	commentstyle=\color{codegray},
	keywordstyle=\color{keywordcolor}\bfseries,
	numberstyle=\tiny\color{codegray},
	stringstyle=\color{codepurple},
	basicstyle=\ttfamily\footnotesize,
	breakatwhitespace=false,
	breaklines=true,
	captionpos=b,
	keepspaces=true,
	numbers=left,
	numbersep=5pt,
	showspaces=false,
	showstringspaces=false,
	showtabs=false,
	tabsize=2,
	language=Java,
	frame=single,
	rulecolor=\color{black!30}
}

\lstset{style=javastyle}

% --- Hinweis-/Beispiel-/Formel-Umgebungen (schlicht, ohne farbige Box) ---
\newenvironment{hinweis}[1][Hinweis]%
	{\par\smallskip\noindent\textbf{#1.}\enspace\ignorespaces}%
	{\par\smallskip}

\newenvironment{beispiel}[1][Beispiel]%
	{\par\smallskip\noindent\textbf{#1:}\par\nobreak\vspace{2pt}}%
	{\par\smallskip}

\newenvironment{formeln}[1][Gegebene Formeln]%
	{\par\smallskip\noindent\textbf{#1:}\par\nobreak\vspace{2pt}}%
	{\par\smallskip}

% ##################################################
% Struktogramm (Nassi-Shneiderman) in reinem TikZ
% ##################################################
\newlength{\nswidth}\setlength{\nswidth}{11.5cm}   % Gesamtbreite
\newlength{\nsind}\setlength{\nsind}{0.55cm}        % Einrueckung je Schleifenebene
% vorberechnete Versatzbreiten fuer Verzweigungen (in Koordinaten nutzbar)
\newlength{\nshalf}\setlength{\nshalf}{\dimexpr\nswidth/2\relax}            % halbe Gesamtbreite
\newlength{\nshalfi}\setlength{\nshalfi}{\dimexpr(\nswidth-\nsind)/2\relax}  % halbe Breite, 1x eingerueckt
\newlength{\nshalfii}\setlength{\nshalfii}{\dimexpr(\nswidth-2\nsind)/2\relax} % halbe Breite, 2x eingerueckt

\tikzset{
	% Basis fuer Anweisungs-/Verzweigungsbloecke
	nb/.style={draw, line width=0.4pt, anchor=north west, inner sep=4pt,
		align=left, font=\footnotesize, fill=white, minimum height=0.8cm},
	% volle Breite
	nbfull/.style={nb, minimum width=\nswidth,
		text width=\dimexpr\nswidth-8pt\relax},
	% eine Ebene eingerueckt (Schleifenrumpf)
	nbi/.style={nb, minimum width=\dimexpr\nswidth-\nsind\relax,
		text width=\dimexpr\nswidth-\nsind-8pt\relax},
	% zwei Ebenen eingerueckt (Schleife in Schleife)
	nbii/.style={nb, minimum width=\dimexpr\nswidth-2\nsind\relax,
		text width=\dimexpr\nswidth-2\nsind-8pt\relax},
	% Verzweigung: halbe Breiten
	nbh/.style={nb, minimum width=\dimexpr\nswidth/2\relax,
		text width=\dimexpr\nswidth/2-8pt\relax},
	nbih/.style={nb, minimum width=\dimexpr(\nswidth-\nsind)/2\relax,
		text width=\dimexpr(\nswidth-\nsind)/2-8pt\relax},
	% Verzweigung halbe Breite, zwei Ebenen eingerueckt
	nbiih/.style={nb, minimum width=\nshalfii,
		text width=\dimexpr\nshalfii-8pt\relax},
	% Kopf einer Verzweigung (das Dreieck), inner sep 0
	nc/.style={draw, line width=0.4pt, anchor=north west, inner sep=0pt,
		minimum height=1.0cm, fill=white},
	ncfull/.style={nc, minimum width=\nswidth},
	nci/.style={nc, minimum width=\dimexpr\nswidth-\nsind\relax},
	ncii/.style={nc, minimum width=\dimexpr\nswidth-2\nsind\relax},
}

% Zeichnet das Dreieck + Beschriftungen in einen bereits gesetzten nc-Knoten.
% #1 Knotenname  #2 Bedingung  #3 ja-Label  #4 nein-Label
\newcommand{\nscond}[4]{%
	\draw[line width=0.4pt] (#1.north west) -- (#1.south);
	\draw[line width=0.4pt] (#1.north east) -- (#1.south);
	\node[font=\scriptsize, anchor=north, inner sep=1pt] at ($(#1.north)+(0,-1.5pt)$) {#2};
	\node[font=\scriptsize, anchor=south west, inner sep=1pt] at ($(#1.south west)+(2.5pt,1.5pt)$) {#3};
	\node[font=\scriptsize, anchor=south east, inner sep=1pt] at ($(#1.south east)+(-2.5pt,1.5pt)$) {#4};
}

% Zeichnet den L-Rahmen einer Schleife (linker + unterer Rand der Einrueckspalte).
% #1 Kopfknoten der Schleife  #2 letzter Rumpfknoten
\newcommand{\nsframe}[2]{%
	\draw[line width=0.4pt] (#1.south west) |- (#2.south west);
}


\title{Probeklausur 5}
\author{Luis Staudt}
\date{}

\begin{document}

% --- Titelblock ---
\begin{center}
	\rule{\textwidth}{0.8pt}\\[0.6em]
	{\LARGE\bfseries Probeklausur 5}\\[0.3em]
	{\large Grundlagen der Programmierung}\\[0.7em]
	\begin{tabular}{r l @{\hspace{1.4cm}} r l}
		\textbf{Studiengänge:}    & PEB \& MVB            & \textbf{Autor:}        & Luis Staudt \\
		\textbf{Semester:}        & Sommersemester 2026   & \textbf{Schwierigkeit:}& schwer \\
		\textbf{Bearbeitungszeit:}& 90 Minuten           & \textbf{Gesamtpunkte:} & 90 \\
	\end{tabular}\\[0.7em]
	\rule{\textwidth}{0.8pt}
\end{center}
\vspace{0.5em}

\noindent\textit{Bearbeiten Sie alle fünf Aufgaben. Achten Sie auf sauberen, kompilierfähigen Java-Quelltext. Hilfsmittel: OpenBook, keine elektronischen Geräte.}

\pagestyle{fancy}
\fancyhf{}
\lhead{\small Probeklausur 5, Grundlagen der Programmierung}
\rhead{\small SoSe 2026 (schwer)}
\cfoot{\thepage}
\renewcommand{\headrulewidth}{0.4pt}

% =============================================
\clearpage
\section*{Aufgabe 1: Struktogramm in Java umsetzen \hfill (24 Punkte)}

Das folgende Struktogramm liest Zahlen ein, klassifiziert sie über ihre Teilersumme (vollkommen / abundant) und sucht anschließend die \emph{erste} Primzahl im Feld.
Beide Teile verwenden eine \textbf{geschachtelte} Schleife.

Setzen Sie das Struktogramm \textbf{vollständig in ein lauffähiges Java-Programm} um.

\begin{center}
\begin{tikzpicture}
	\node[nbfull] (a1) at (0,0) {lies $n$ ein; lege Feld \texttt{werte} mit $n$ Plätzen an};
	\node[nbfull] (h1) at (a1.south west) {\textbf{für} $i = 0,\,1,\,\ldots,\,n-1$};
	\node[nbi]    (b1) at ($(h1.south west)+(\nsind,0)$) {lies \texttt{werte[i]} ein};
	\nsframe{h1}{b1}

	% Schleife 2: Teilersumme + Klassifikation
	\node[nbfull] (h2) at ($(b1.south west)+(-\nsind,0)$) {\textbf{für} $i = 0,\,1,\,\ldots,\,n-1$};
	\node[nbi]    (f0) at ($(h2.south west)+(\nsind,0)$) {\texttt{tsum = 0;}};
	\node[nbi]    (h2i) at (f0.south west) {\textbf{für} $t = 1,\,\ldots,\,\texttt{werte[i]}/2$};
	\node[ncii]   (ic) at ($(h2i.south west)+(\nsind,0)$) {};
	\nscond{ic}{\texttt{werte[i] \% t == 0}}{ja}{nein}
	\node[nbiih]  (icl) at (ic.south west) {\texttt{tsum = tsum + t;}};
	\node[nbiih]  (icr) at ($(ic.south west)+(\nshalfii,0)$) {\textit{(nichts)}};
	\nsframe{h2i}{icl}
	\node[nci]    (c1) at ($(icl.south west)+(-\nsind,0)$) {};
	\nscond{c1}{\texttt{tsum == werte[i]}}{ja}{nein}
	\node[nbih]   (c1l) at (c1.south west) {drucke \texttt{werte[i]}: vollkommen};
	\node[nbih]   (c1r) at ($(c1.south west)+(\nshalfi,0)$) {\textit{(nichts)}};
	\node[nci]    (c2) at (c1l.south west) {};
	\nscond{c2}{\texttt{tsum > werte[i]}}{ja}{nein}
	\node[nbih]   (c2l) at (c2.south west) {\texttt{anzAbund = anzAbund + 1;}};
	\node[nbih]   (c2r) at ($(c2.south west)+(\nshalfi,0)$) {\textit{(nichts)}};
	\nsframe{h2}{c2l}

	\node[nbfull] (a3) at ($(c2l.south west)+(-\nsind,0)$) {\texttt{i = 0;}\quad \texttt{primGefunden = false;}};

	% Schleife 3: erste Primzahl
	\node[nbfull] (h3) at (a3.south west) {\textbf{solange} (\texttt{nicht primGefunden}) \textbf{und} (\texttt{i < n})};
	\node[nbi]    (pf) at ($(h3.south west)+(\nsind,0)$) {\texttt{istPrim = (werte[i] >= 2);}};
	\node[nbi]    (h3i) at (pf.south west) {\textbf{für} $t = 2,\,\ldots,\,\texttt{werte[i]}-1$};
	\node[ncii]   (pic) at ($(h3i.south west)+(\nsind,0)$) {};
	\nscond{pic}{\texttt{werte[i] \% t == 0}}{ja}{nein}
	\node[nbiih]  (picl) at (pic.south west) {\texttt{istPrim = false;}};
	\node[nbiih]  (picr) at ($(pic.south west)+(\nshalfii,0)$) {\textit{(nichts)}};
	\nsframe{h3i}{picl}
	\node[nci]    (pc) at ($(picl.south west)+(-\nsind,0)$) {};
	\nscond{pc}{\texttt{istPrim == true}}{ja}{nein}
	\node[nbih]   (pcl) at (pc.south west) {\texttt{primGefunden = true; primPos = i;}};
	\node[nbih]   (pcr) at ($(pc.south west)+(\nshalfi,0)$) {\textit{(nichts)}};
	\node[nbi]    (pinc) at (pcl.south west) {\texttt{i = i + 1;}};
	\nsframe{h3}{pinc}

	% Ausgabe
	\node[ncfull] (e1) at ($(pinc.south west)+(-\nsind,0)$) {};
	\nscond{e1}{\texttt{primGefunden == true}}{ja}{nein}
	\node[nbh]    (e1l) at (e1.south west) {drucke \texttt{primPos}};
	\node[nbh]    (e1r) at ($(e1.south west)+(\nshalf,0)$) {drucke: keine Primzahl};
\end{tikzpicture}
\end{center}

% =============================================
\clearpage
\section*{Aufgabe 2: Schiefer Wurf, Größe auswählen \hfill (14 Punkte)}

Ein Körper wird mit der Anfangsgeschwindigkeit $v_0$ unter dem Winkel~$\alpha$ (in Grad) geworfen ($g = 9{,}81$).

\begin{formeln}
\[
	\text{(1) Horizontalweg:}\;\; x = v_0 \cos\alpha \cdot t
	\qquad
	\text{(2) Höhe:}\;\; y = v_0 \sin\alpha \cdot t - \tfrac{1}{2} g\,t^2
\]
\[
	\text{(3) Steigzeit:}\;\; t_s = \frac{v_0 \sin\alpha}{g}
	\qquad
	\text{(4) Wurfweite:}\;\; w = \frac{v_0^{\,2}\,\sin(2\alpha)}{g}
\]
\end{formeln}

Schreiben Sie ein Programm, das die gesuchte Größe per Zahl (\texttt{1} bis \texttt{4}) wählen lässt, die benötigten Werte abfragt und mit der \textbf{richtigen} Formel rechnet.
Behandeln Sie die Sonderfälle: bei $\alpha = 90^\circ$ ist der Horizontalweg~$x = 0$ (und die Wurfweite~$w = 0$), bei $\alpha = 0^\circ$ ist die Steigzeit $t_s = 0$.

\begin{hinweis}[Winkel im Bogenmaß]
	Die Java-Funktionen \texttt{Math.sin} und \texttt{Math.cos} erwarten das Bogenmaß.
	Rechnen Sie den Winkel mit \texttt{Math.toRadians(alpha)} bzw.\ $\alpha \cdot \pi/180$ um.
\end{hinweis}

% =============================================
\clearpage
\section*{Aufgabe 3: Klasse \texttt{Dreieck} \hfill (20 Punkte)}

Gegeben ist die Klasse \texttt{Punkt} (Attribute \texttt{xKoordinate}, \texttt{yKoordinate} als \texttt{double}).
Implementieren Sie eine Klasse \texttt{Dreieck} aus drei \texttt{Punkt}-Objekten.
Das Gerüst ist vorgegeben. Ergänzen Sie:

\begin{itemize}[nosep]
	\item einen \textbf{Konstruktor} mit drei \texttt{Punkt}-Parametern.
	\item \texttt{berechneFlaeche()}:
		$A = \tfrac{1}{2}\,\bigl|\,\Delta x_{12}\cdot\Delta y_{13} - \Delta x_{13}\cdot\Delta y_{12}\,\bigr|$,
		wobei $\Delta x_{12} = x_2 - x_1$ usw.
	\item \texttt{berechneUmfang()}: Summe der drei Seitenlängen (\emph{verwenden Sie}\\ \texttt{berechneAbstandZuPunkt}).
	\item \texttt{istRechtwinklig()}: gibt \texttt{true} zurück, wenn das Dreieck (über den Satz des Pythagoras) rechtwinklig ist.
		Wegen Rundungsfehlern ist eine \textbf{Epsilon-Umgebung} zu verwenden.
\end{itemize}

\begin{lstlisting}
public class Dreieck {
    public Punkt p1, p2, p3;   // Daten
    // Konstruktor und Methoden ergaenzen
}
\end{lstlisting}

\begin{lstlisting}
Punkt a = new Punkt();  a.xKoordinate = 0;  a.yKoordinate = 0;
Punkt b = new Punkt();  b.xKoordinate = 4;  b.yKoordinate = 0;
Punkt c = new Punkt();  c.xKoordinate = 0;  c.yKoordinate = 3;

Dreieck d = new Dreieck(a, b, c);
System.out.println("Flaeche: " + d.berechneFlaeche());  // 6.0
System.out.println("Umfang:  " + d.berechneUmfang());   // 12.0
if (d.istRechtwinklig())
    System.out.println("Das Dreieck ist rechtwinklig");
\end{lstlisting}

% =============================================
\clearpage
\section*{Aufgabe 4: Fachwerk, Knotenverschiebungen \hfill (16 Punkte)}

Gegeben ist die aus Vorlesung und Übung bekannte Datei \texttt{fachwerk.jar}.
Nach der Berechnung sollen für ein gewähltes Fachwerk zwei Auswertungen der Knotenverschiebungen programmiert werden:
\begin{itemize}[nosep]
	\item Geben Sie \textbf{alle Knoten} aus, die sich \textbf{rein vertikal} verschieben (Verschiebung in x-Richtung praktisch null, in y-Richtung von null verschieden). Existiert kein solcher Knoten, ist eine entsprechende Meldung auszugeben.
	\item Geben Sie zusätzlich den Knoten mit der \textbf{größten Gesamtverschiebung} samt deren Wert aus, wobei $v = \sqrt{v_x^2 + v_y^2}$.
\end{itemize}

\begin{center}
\begin{tikzpicture}[>=Stealth, scale=1.0,
	knoten/.style={circle, fill=black, inner sep=1.6pt},
	stab/.style={thick, gray!70!black}]
	\coordinate (A) at (0,0); \coordinate (B) at (2.5,0); \coordinate (C) at (5,0);
	\coordinate (D) at (1.25,2.0); \coordinate (E) at (3.75,2.0);
	\draw[stab] (A)--(B); \draw[stab] (B)--(C);
	\draw[stab] (A)--(D); \draw[stab] (D)--(B); \draw[stab] (D)--(E);
	\draw[stab] (E)--(B); \draw[stab] (E)--(C);
	\foreach \p in {A,B,C,D,E} \node[knoten] at (\p) {};
	\draw[->, thick] (B)--++(0,-0.9);
	\node[below, font=\footnotesize] at ($(B)+(0,-0.9)$) {Last};
\end{tikzpicture}\\[-0.3em]
{\footnotesize\itshape Schematische Darstellung eines Fachwerks.}
\end{center}

\begin{hinweis}[Epsilon-Umgebung beachten]
	Wegen Gleitkomma-Rundung ist \glqq null\grqq{} über eine Epsilon-Umgebung zu prüfen:
	\[ |v_x| < \varepsilon \;\text{ und }\; |v_y| \ge \varepsilon. \]
	Benötigte Methoden: \texttt{gibAnzahlKnoten()}, \texttt{gibKnoten(int i)},\\ \texttt{gibVerschiebungX()}, \texttt{gibVerschiebungY()}.
\end{hinweis}

Ergänzen Sie den folgenden Quelltext zu einem vollständigen, lauffähigen Programm:

\begin{lstlisting}
public class Fachwerkaufgabe {
    public static void main(String[] args) {
        System.out.println("Welches Fachwerk wollen Sie rechnen (0-4)?");
        int wahl = Tastatureingabe.leseInt();
        Fachwerk fw = new Fachwerk(wahl);
        fw.berechneGroessen();
        double epsilon = 1e-6;
        // hier erweitern: rein vertikal verschobene Knoten ausgeben sowie
        // den Knoten mit der groessten Gesamtverschiebung bestimmen
    }
}
\end{lstlisting}

\begin{beispiel}[Mögliche Ausgabe]
\begin{lstlisting}[language={}]
Knoten 2 verschiebt sich rein vertikal
Knoten 4 verschiebt sich rein vertikal
Groesste Gesamtverschiebung: Knoten 3 mit 0.0123
\end{lstlisting}
\end{beispiel}

% =============================================
\clearpage
\section*{Aufgabe 5: Wurfbahn in eine Datei schreiben \hfill (16 Punkte)}

Ein Körper wird unter dem Winkel~$\alpha$ mit $v_0$ geworfen ($g = 9{,}81$). Für die Bahnkurve gilt
\[ x(t) = v_0 \cos\alpha \cdot t \qquad y(t) = v_0 \sin\alpha \cdot t - \tfrac{1}{2} g\,t^2. \]

Schreiben Sie ein Programm, das ab $t = 0~\text{s}$ in Schritten von $0{,}1~\text{s}$ die Werte $t$, $x(t)$ und $y(t)$ in die Datei \texttt{bahn.txt} schreibt. \textbf{Brechen Sie ab, sobald $y(t) < 0$} wird (Aufprall).
Hängen Sie am Ende eine Zeile mit der erreichten \textbf{Wurfweite} und der \textbf{Flugzeit} an.

\begin{hinweis}[Hinweise]
	Winkel mit \texttt{Math.toRadians(alpha)} umrechnen.
	Datei: \lstinline|new FileWriter("bahn.txt")|, je Zeile \texttt{write(...)} mit \lstinline|"\n"|, am Ende \texttt{close()}.
	Nutzen Sie \lstinline|"\t"| für die Spaltenabstände.
\end{hinweis}

\begin{beispiel}[Erwarteter Dateiinhalt (Auszug, $v_0=20$, $\alpha=45^\circ$)]
\begin{lstlisting}[language={}]
t [s]   x [m]    y [m]
0.0     0.0      0.0
0.1     1.414    1.365
...
2.8     39.6     0.179
Aufprall bei t = 2.9 s, Wurfweite = 41.01 m
\end{lstlisting}
\end{beispiel}

\end{document}
